% !TeX program = pdflatex
% A proof of the Elizalde--Luo conjecture for nonnesting permutations avoiding 1132 and 3312
% Build: pdflatex note.tex (run twice). See BUILD.md.
\documentclass[11pt]{article}

\usepackage[T1]{fontenc}
\usepackage{lmodern}
\usepackage[margin=1in]{geometry}
\usepackage{amsmath,amssymb,amsthm}
\usepackage{booktabs}
\usepackage{array}
\usepackage{microtype}
\usepackage[hidelinks]{hyperref}
\usepackage{xurl}

\newtheorem{theorem}{Theorem}[section]
\newtheorem{lemma}[theorem]{Lemma}
\newtheorem{proposition}[theorem]{Proposition}
\newtheorem{corollary}[theorem]{Corollary}
\theoremstyle{definition}
\newtheorem{definition}[theorem]{Definition}
\newtheorem{example}[theorem]{Example}
\theoremstyle{remark}
\newtheorem{remark}[theorem]{Remark}

\newcommand{\nn}{[n]_2}
\DeclareMathOperator{\fst}{fst}
\DeclareMathOperator{\snd}{snd}
\newcommand{\Av}{\operatorname{Av}}
\newcommand{\sU}{\mathsf{U}}
\newcommand{\sD}{\mathsf{D}}
\newcommand{\sL}{\mathsf{L}}
\newcommand{\sH}{\mathsf{H}}
\newcommand{\sA}{\mathsf{A}}
\newcommand{\sB}{\mathsf{B}}
\newcommand{\sC}{\mathsf{C}}
\newcommand{\Cat}{C}
\newcommand{\Wn}{W_n}
\newcommand{\open}{\operatorname{open}}

\title{A proof of the Elizalde--Luo conjecture for nonnesting\\ permutations avoiding $1132$ and $3312$}
\author{John Erlbacher\thanks{Independent Researcher. Email: \texttt{erlbacher.research@gmail.com}. ORCID: 0009-0003-6851-4139.}}
\date{June 2026}

\begin{document}
\maketitle

\begin{abstract}
Nonnesting permutations are permutations of the multiset
$\{1,1,2,2,\dots,n,n\}$ that avoid the patterns $1221$ and $2112$; there are
$n!\,\Cat_n$ of them, where $\Cat_n$ is the $n$th Catalan number. In the
concluding section of their systematic study of pattern avoidance in
nonnesting permutations [DMTCS 27:1 (2025), \#13], Elizalde and Luo
conjectured, on the basis of data for $n\le 8$, that the number of nonnesting
permutations of $\{1,1,\dots,n,n\}$ avoiding both $1132$ and $3312$ equals
$3^n-3\cdot 2^{n-1}+1$ (OEIS A168583, shifted). We prove this conjecture. The
proof is elementary and self-contained: nonnesting words factor as a Dyck
shape carrying a label permutation; avoidance of $\{1132,3312\}$ forces the
labels to form a prefix-interval permutation, encoded by a sign word in
$\{\sL,\sH\}^{n-1}$, and reduces to a two-coloring condition on the crossing
arcs of the shape. The shapes admitting a valid coloring fall into three
explicit families, and summing the number of colorings gives the formula. We
also give a second, bijective proof of the final count: an explicit
letter-by-letter bijection between the avoiders and a ternary language
$\Wn\subseteq\{\sA,\sB,\sC\}^n$ whose cardinality is
$3^n-3\cdot2^{n-1}+1$ by inspection. Every lemma has been verified
mechanically far beyond the conjecture's original range, including an
exhaustive word-level check of the structural characterization over all
$57{,}657{,}600$ nonnesting words for $n=8$.
\end{abstract}

\section{Introduction}\label{sec:intro}

A \emph{Stirling permutation} \cite{GesselStanley1978} is a permutation of
the multiset $\nn=\{1,1,2,2,\dots,n,n\}$ in which the letters between the two
copies of any value are larger than that value. Two natural relaxations have
attracted attention in recent years: \emph{noncrossing} (or
\emph{quasi-Stirling}) permutations, the permutations of $\nn$ with no
subsequence $abab$ with $a\ne b$, introduced by Archer, Gregory, Pennington
and Slayden \cite{Archer2019} and studied further in
\cite{Elizalde2021}; and \emph{nonnesting} permutations, those with no
subsequence $abba$ with $a\ne b$, introduced by Elizalde \cite{Elizalde2024}.
There are exactly $n!\,\Cat_n$ nonnesting permutations of $\nn$, where
$\Cat_n=\frac{1}{n+1}\binom{2n}{n}$ \cite{Elizalde2024,ElizaldeLuo}; we
rederive this in Lemma~\ref{lem:fifo} below.

Elizalde and Luo \cite{ElizaldeLuo}, extending Luo's thesis \cite{Luo2024},
carried out a systematic study of pattern avoidance in nonnesting
permutations: they enumerated the nonnesting permutations avoiding each set
of two or more patterns of length $3$, as well as those avoiding several sets
of patterns of length $4$. (For single patterns of length $3$ the noncrossing
and nonnesting problems were subsequently advanced by Archer and Laudone
\cite{ArcherLaudone}.) In their concluding section, Elizalde and Luo list a
table of conjectured enumeration formulas for five further sets of length-$4$
patterns \cite[Table~4]{ElizaldeLuo}, each checked by computer for
$n\le 8$. The present paper proves the conjecture for the set
$\{1132,3312\}$.

Throughout, a \emph{word} (or \emph{permutation of $\nn$}) is a sequence
$w=w_1w_2\cdots w_{2n}$ over $\{1,\dots,n\}$ using each letter exactly twice,
and containment of a pattern is understood in the sense of
\cite{ElizaldeLuo}, recalled in Section~\ref{sec:prelim}: equalities
among pattern letters must be matched by equalities among word letters, and
strict inequalities by strict inequalities. A word is \emph{nonnesting} if it
avoids $1221$ and $2112$. Following \cite{ElizaldeLuo}, $\mathcal C_n$
denotes the set of nonnesting permutations of $\nn$, and $c_n(\Lambda)$ the
number of elements of $\mathcal C_n$ avoiding every pattern in a set
$\Lambda$.

\begin{theorem}[Elizalde--Luo conjecture for $\{1132,3312\}$]\label{thm:main-intro}
For every $n\ge 1$,
\[
  c_n(\{1132,3312\}) \;=\; 3^n-3\cdot 2^{n-1}+1 .
\]
\end{theorem}

The conjecture appears as the row
``$\{1132,3312\}$ \;---\; $3^n-3\cdot2^{n-1}+1$ \;---\; A168583''
of \cite[Table~4]{ElizaldeLuo}, whose caption reads \emph{``Some conjectures
on the enumeration of nonnesting permutations avoiding other patterns''} and
whose preamble states that \emph{``All the conjectures have been checked for
$n$ up to $8$.''} The sequence
$1,4,16,58,196,634,1996,6178,\dots$ given by the formula is the shift
$a(n{+}2)$ of OEIS entry A168583 \cite{OEIS}, whose $m$th term ($m\ge3$, the
entry's own indexing) counts the partitions of the multiset
$\{1,1,2,3,\dots,m-1\}$ into exactly three nonempty parts; as a byproduct,
Theorem~\ref{thm:main-intro} shows that the avoiders of $\nn$ are
equinumerous with the partitions of $\{1,1,2,3,\dots,n+1\}$ into exactly
three nonempty parts (an equality of cardinalities via the closed form; we
exhibit no direct bijection). To our
knowledge the conjecture was open before this work: it is stated as a
conjecture in the published version of \cite{ElizaldeLuo} (October 2025)
and is unchanged in the latest arXiv revision (v6, January 2026, retrieved
June 12, 2026), and we are aware of no announcement of a proof (literature
and web searches last performed June 12, 2026).

\subsection*{Outline of the proof}

All of the work happens in an explicit, finite combinatorial setting; no
generating-function or analytic machinery is needed.

\begin{enumerate}
\item \textbf{FIFO normal form (Section~\ref{sec:prelim}).} A word is
nonnesting if and only if, for every $i$, its $i$th closer (second occurrence
of a value, in position order) carries the same letter as its $i$th opener.
Consequently nonnesting words correspond bijectively to pairs $(s,p)$ where
$s$ is a Dyck word of semilength $n$ (the \emph{shape}) and $p\in S_n$ is the
sequence of labels read along the openers. We also unwind the patterns
$1132$ and $3312$ into a positional criterion: $w$ contains one of them iff
some letter lies strictly between an earlier-closed value $v$ and a letter
positioned in between (Corollary~\ref{cor:combined}).
\item \textbf{Sign words (Section~\ref{sec:sign}).} If $(s,p)$ avoids
$\{1132,3312\}$, the label permutation $p$ must be \emph{prefix-interval}:
every prefix of $p$ is an interval of integers. Such $p$ are classically the
$\{132,312\}$-avoiding permutations; they are encoded by a \emph{sign word}
$\varepsilon\in\{\sL,\sH\}^{n-1}$ recording whether each new label is the
current minimum or maximum.
\item \textbf{The characterization (Section~\ref{sec:thmA}).}
Theorem~\ref{thm:A}: $(s,p)$ is an avoider iff $p$ is prefix-interval and its
sign word satisfies
\[
 (C)\colon\qquad \varepsilon_K\neq\varepsilon_J\quad
 \text{for every pair of arcs } K<J \text{ with } q_1<o_J<q_K,
\]
where $o_i,q_i$ are the positions of the $i$th opener and closer. The
pairs in $(C)$ are crossing pairs of arcs whose later arc opens after the
first closer.
\item \textbf{Shape classification (Section~\ref{sec:shapes}).} The Dyck
words $s$ for which $(C)$ is satisfiable form three explicit families
(Lemma~\ref{lem:classification}), parametrized by the length $a$ of the first
ascent, an optional ``bridge'' position $i$, and a vector of tail heights
$\delta\in\{0,1\}^{n-a-1}$; on these shapes the valid sign words are exactly
the free choices of $|F(s)|$ coordinates, so their number is $2^{|F(s)|}$
(Corollary~\ref{cor:free}).
\item \textbf{Two endgames.} Summing $2^{|F(s)|}$ over the three families
gives $3^n-3\cdot2^{n-1}+1$ (Section~\ref{sec:count}). Independently,
Section~\ref{sec:bijection} packages the classification data into a single
ternary string and exhibits an explicit bijection, with an explicit inverse,
between the avoiders and the language
\[
 \Wn \;=\; \{\sA,\sB,\sC\}^n\setminus
 \bigl(\sB\{\sA,\sB\}^{n-1}\,\cup\,
 \sC\{\sA,\sB\}^*\sC^*\bigr),
\]
whose cardinality is $3^n-2^{n-1}-(2^n-1)$ by a two-line count. The two
endgames are logically independent; either one completes the proof.
\end{enumerate}

Beyond the proofs, every lemma in this paper has been verified mechanically
on ranges far exceeding the $n\le 8$ data that motivated the conjecture,
including: an exhaustive word-level verification of
Theorem~\ref{thm:A} for every one of the $2{,}162{,}160$ nonnesting words at
$n=7$ and every one of the $57{,}657{,}600$ at $n=8$; shape-level
verification of the classification for all $742{,}900$ Dyck shapes of
semilength $13$; and an end-to-end check of the bijection on all
ground-truth avoiders up to $n=10$. Appendix~\ref{app:verification} gives the
full table of checks and how to reproduce the small-$n$ layer from the
literal definitions. Appendix~\ref{app:methods} discloses how this paper was
produced (AI-generated mathematics with adversarial machine refereeing) and
what was and was not verified by machine.

\section{Preliminaries: conventions and the FIFO normal form}\label{sec:prelim}

\subsection{Conventions}\label{sec:conventions}

We work with the definitions of \cite{ElizaldeLuo}, which we now recall. Let $[n]=\{1,\dots,n\}$ and let $\nn=\{1,1,2,2,\dots,n,n\}$ be the
multiset with two copies of each element of $[n]$. A \emph{permutation of
$\nn$}, here called a \emph{word}, is a sequence $w=w_1\cdots w_{2n}$ over
$[n]$ in which every value occurs exactly twice. Positions are $1$-based.

Given a word $w$ and a pattern $\sigma=\sigma_1\cdots\sigma_k$ over the
positive integers, $w$ \emph{contains} $\sigma$ if there exist indices
$1\le i_1<\cdots<i_k\le 2n$ such that for all $r,s\in[k]$,
\[
  w_{i_r}<w_{i_s} \iff \sigma_r<\sigma_s
  \qquad\text{and}\qquad
  w_{i_r}=w_{i_s} \iff \sigma_r=\sigma_s .
\]
Otherwise $w$ \emph{avoids} $\sigma$. Note both conditions are
biconditional: equal pattern letters must map to equal word letters and
conversely.\footnote{Before any proof attempt, these conventions were pinned
in the project's definitions file as verbatim quotations of the published
LaTeX source of \cite{ElizaldeLuo} (the statement above adapts only the
notation), and our computational implementations were validated against the
literal definition (Appendix~\ref{app:verification}).}

A word is \emph{nonnesting} if it avoids $1221$ and $2112$. For a value $v$
write $\fst(v)<\snd(v)$ for the positions of its two occurrences; the
\emph{arc} of $v$ is the pair $(\fst(v),\snd(v))$. A position is an
\emph{opener} if it is the first occurrence of its value, and a
\emph{closer} otherwise. Unwinding the containment convention (equal pattern
letters force equal values, and each value has exactly two copies), $w$ is
nonnesting if and only if there are no two values $u\neq v$ with
\[
  \fst(u)<\fst(v)<\snd(v)<\snd(u),
\]
i.e.\ no arc nests inside another.

A nonnesting word avoiding both $1132$ and $3312$ is called an
\emph{avoider}, and the number of avoiders of $\nn$ is denoted by
$c_n=c_n(\{1132,3312\})$.
Avoidance of $\{1132,3312\}$ is \emph{not} invariant under relabeling of the
values, and the count is over all labeled words, with no quotient
taken.\footnote{For instance, at $n=3$ there are $16$ avoiders but $3$
avoiders whose first occurrences of $1,2,3$ appear in increasing order;
$3\cdot 3!=18\neq 16$, so the familiar ``canonical labeling'' shortcut is
invalid for this problem.}

We say $x$ is \emph{strictly between} $y$ and $z$ if $y<x<z$ or $z<x<y$.
Integer intervals $\{u..v\}=\{j\in\mathbb Z: u\le j\le v\}$ are understood to
be empty when $u>v$.

\subsection{Nonnesting words as labeled Dyck shapes}

Mark each opener of a word $w$ with $\sU$ and each closer with $\sD$; the
resulting word $s=s(w)\in\{\sU,\sD\}^{2n}$ is the \emph{shape} of $w$. Since
distinct closers in a prefix are preceded by their own distinct openers in
the same prefix, every prefix of $s$ has at least as many $\sU$'s as
$\sD$'s; thus $s$ is a \emph{Dyck word} of semilength $n$. Write
$o_1<\cdots<o_n$ for the positions of the $\sU$'s (\emph{openers}) and
$q_1<\cdots<q_n$ for those of the $\sD$'s (\emph{closers}). In any Dyck word,
\begin{equation}\label{eq:oq}
  o_i<q_i \qquad\text{for all } i\in[n],
\end{equation}
because the prefix ending at $q_i$ contains $i$ letters $\sD$, hence at
least $i$ letters $\sU$.

For a Dyck word $s$ and a permutation $p=(p_1,\dots,p_n)$ of $[n]$, let
$\mathrm{wd}(s,p)$ be the word with the letter $p_i$ at positions $o_i$ and
$q_i$ for each $i$.

\begin{lemma}[FIFO normal form]\label{lem:fifo}
A word $w$ over $\nn$ is nonnesting if and only if for every $i\in[n]$ the
letter at the $i$th opener of $w$ equals the letter at the $i$th closer.
Consequently, the map $w\mapsto (s(w),p(w))$, where $p_i$ is the letter at
the $i$th opener, is a bijection from the nonnesting words of $\nn$ onto the
pairs (Dyck word $s$ of semilength $n$, permutation $p\in S_n$), with
inverse $(s,p)\mapsto\mathrm{wd}(s,p)$. In particular
$|\mathcal C_n| = n!\,\Cat_n$.
\end{lemma}

\begin{proof}
The arcs of $w$ form a perfect matching of $[2n]$ pairing each opener with a
closer; let $\sigma\in S_n$ be the permutation such that the value opened at
$o_i$ has its second occurrence at $q_{\sigma(i)}$.

First, $\sigma=\mathrm{id}$ if and only if the letter condition holds.
Indeed, if $\sigma=\mathrm{id}$ then the value at $q_i$ is the one opened at
$o_i$. Conversely, suppose $w_{o_i}=w_{q_i}=:v$ for every $i$. Since $v$
occurs exactly twice, and $o_i\neq q_i$ are an opener and a closer carrying
$v$, the positions $o_i,q_i$ are precisely the two occurrences of $v$; hence
$\sigma(i)=i$.

($\Leftarrow$) Assume $\sigma=\mathrm{id}$ and suppose two values nest:
$\fst(u)<\fst(v)$ and $\snd(v)<\snd(u)$. The first inequality says $u$ is
opened at an earlier opener, $u=p_i$ and $v=p_j$ with $i<j$; but then
$\snd(u)=q_i<q_j=\snd(v)$, a contradiction. So $w$ is nonnesting.

($\Rightarrow$) Assume $\sigma\neq\mathrm{id}$ and pick $i<j$ with
$\sigma(i)>\sigma(j)$. The value opened at $o_j$ closes at $q_{\sigma(j)}$,
and since the two positions of any one value satisfy $\fst<\snd$, we have
$o_j<q_{\sigma(j)}$; hence
\[
  o_i \,<\, o_j \,<\, q_{\sigma(j)} \,<\, q_{\sigma(i)},
\]
so the value opened at $o_j$ nests inside the value opened at $o_i$: $w$ is
not nonnesting.

For the bijection: given nonnesting $w$, the above gives
$w=\mathrm{wd}(s(w),p(w))$. Conversely, let $s$ be a Dyck word and
$p\in S_n$, and consider $\mathrm{wd}(s,p)$. The position $o_i$ is the first
occurrence of $p_i$: any earlier position carries some $p_j$ with $j<i$ (an
earlier $\sU$ is the $j$th with $j<i$; an earlier $\sD$ is the $j$th $\sD$,
which by the Dyck property is preceded by at least $j$ $\sU$'s, so $j\le
i-1$). Hence the arc of $p_i$ is $(o_i,q_i)$ by \eqref{eq:oq}; arcs sorted
simultaneously by opener and closer never nest, so $\mathrm{wd}(s,p)$ is
nonnesting, and its shape and label data are $(s,p)$. The two constructions
are mutually inverse, and the count $|\mathcal C_n|=\Cat_n\cdot n!$ follows.
\end{proof}

From now on we identify a nonnesting word with its pair $(s,p)$ and call
$i\in[n]$ the \emph{arc} with endpoints $(o_i,q_i)$ and label $p_i$. For
$i<j$, arcs $i,j$ are \emph{crossing} if $o_j<q_i$ and \emph{disjoint} if
$q_i<o_j$; by Lemma~\ref{lem:fifo} nesting cannot occur.

\begin{remark}[arc support]\label{rem:support}
Arc $i$ has no letters at positions $>q_i$, and every position $z>q_i$
carries a letter of an arc $J>i$: if $z=o_J$ then $q_J>o_J>q_i$, and if
$z=q_J$ then $q_J>q_i$; in either case $q_J>q_i$, which since closers are
sorted forces $J>i$.
\end{remark}

\subsection{Unwinding the patterns \texorpdfstring{$1132$ and $3312$}{1132 and 3312}}

\begin{lemma}\label{lem:positional}
Let $w$ be any word over $\nn$. Then:
\begin{enumerate}
\item $w$ contains $1132$ $\iff$ there are a value $a$ and positions
  $\snd(a)<k<l$ with $a<w_l<w_k$;
\item $w$ contains $3312$ $\iff$ there are a value $c$ and positions
  $\snd(c)<k<l$ with $w_k<w_l<c$.
\end{enumerate}
\end{lemma}

\begin{proof}
(1, $\Rightarrow$) Let $i_1<i_2<i_3<i_4$ be an occurrence of $1132$. The
pattern letters at $i_1,i_2$ are equal ($1,1$), so $w_{i_1}=w_{i_2}=:a$;
since $a$ has exactly two copies, $\{i_1,i_2\}=\{\fst(a),\snd(a)\}$ and
$i_2=\snd(a)$. Take $k=i_3$, $l=i_4$. The biconditional convention forces
$w_{i_3}>w_{i_4}$ (pattern $3>2$) and $w_{i_4}>w_{i_1}=a$ (pattern $2>1$),
i.e.\ $a<w_l<w_k$.

(1, $\Leftarrow$) The quadruple of positions
$(\fst(a),\snd(a),k,l)$ is increasing, and its letters $(a,a,w_k,w_l)$
realize exactly the order and equality relations of $(1,1,3,2)$: the first
two letters are equal and smaller than both others, and $w_k>w_l$, with all
stated inequalities strict. So it is an occurrence of $1132$.

(2) is the same argument with all inequalities above reversed.
\end{proof}

\begin{corollary}[combined criterion]\label{cor:combined}
A word $w$ over $\nn$ avoids $\{1132,3312\}$ if and only if there is
\emph{no} triple $(v,x,y)$ with $v$ a value and $\snd(v)<x<y\le 2n$ such
that $w_y$ is strictly between $v$ and $w_x$.
\end{corollary}

\begin{proof}
``$w_y$ strictly between $v$ and $w_x$'' means $v<w_y<w_x$ (an occurrence of
$1132$ with $a=v$, by Lemma~\ref{lem:positional}(1)) or $w_x<w_y<v$ (an
occurrence of $3312$ with $c=v$, by Lemma~\ref{lem:positional}(2)).
\end{proof}

\section{The sign word}\label{sec:sign}

For a permutation $p\in S_n$, an occurrence of the (classical) pattern $132$
is a triple of indices $i<j<k$ with $p_i<p_k<p_j$, and an occurrence of
$312$ is $i<j<k$ with $p_j<p_k<p_i$. Having either one means: $p_k$ is
strictly between $p_i$ and $p_j$ for some pair $i,j<k$.

\begin{lemma}\label{lem:labels}
Let $w=(s,p)$ be nonnesting. If $p$ contains $132$ or $312$, then $w$
contains $1132$ or $3312$. Hence the label permutation of an avoider avoids
$\{132,312\}$.
\end{lemma}

\begin{proof}
Suppose $i<j<k$ with $p_k$ strictly between $p_i$ and $p_j$. The positions
$q_i<q_j<q_k$ carry the letters $p_i,p_j,p_k$, and $q_i=\snd(p_i)$ by
Lemma~\ref{lem:fifo}. Apply Corollary~\ref{cor:combined} with
$(v,x,y)=(p_i,q_j,q_k)$: $w_y=p_k$ is strictly between $v=p_i$ and
$w_x=p_j$, so $w$ contains a forbidden pattern.
\end{proof}

Write $\Av_n(132,312)$ for the set of $p\in S_n$ avoiding both patterns.

\begin{lemma}[structure of $\Av(132,312)$]\label{lem:signword}
For $p\in S_n$ the following are equivalent:
\begin{enumerate}
\item $p$ avoids $\{132,312\}$;
\item every $p_k$ is the minimum or the maximum of $\{p_1,\dots,p_k\}$;
\item every prefix set $\{p_1,\dots,p_k\}$ is an interval of integers
  (``$p$ is prefix-interval'').
\end{enumerate}
Moreover, the map
$p\mapsto\varepsilon(p)=(\varepsilon_2,\dots,\varepsilon_n)\in\{\sL,\sH\}^{n-1}$,
where $\varepsilon_k=\sH$ if $p_k=\max\{p_1,\dots,p_k\}$ and
$\varepsilon_k=\sL$ if $p_k=\min\{p_1,\dots,p_k\}$, is a bijection from
$\Av_n(132,312)$ onto $\{\sL,\sH\}^{n-1}$, and for all $i<j$,
\begin{equation}\label{eq:star}
  p_i<p_j \iff \varepsilon_j=\sH
  \qquad\text{and}\qquad
  p_i>p_j \iff \varepsilon_j=\sL.
\end{equation}
We call $\varepsilon(p)$ the \emph{sign word} of $p$, and write
$\bar{x}$ for the letter of $\{\sL,\sH\}$ other than $x$.
\end{lemma}

\begin{proof}
(1)$\Rightarrow$(2): if $p_k$ is not extreme in its prefix, there are
$u,v<k$ with $p_u<p_k<p_v$; if $u<v$ then $(u,v,k)$ is an occurrence of
$132$, and if $v<u$ then $(v,u,k)$ is an occurrence of $312$.

(2)$\Rightarrow$(3): if some prefix $\{p_1,\dots,p_k\}$ is not an interval,
it omits a value $g$ strictly between its minimum and maximum. Since $p$ is
a permutation of $[n]$, $g=p_l$ for some $l>k$; then $p_l$ is neither the
minimum nor the maximum of $\{p_1,\dots,p_l\}$, because the prefix of index
$k$ already contains values on both sides of $g$. This contradicts (2).

(3)$\Rightarrow$(1): under (3), each $p_k$ extends the interval
$\{p_1,\dots,p_{k-1}\}$ by one value at the bottom or the top, hence is
never strictly between two earlier values; so no occurrence of $132$ or
$312$ exists.

Bijection: given (3), $\varepsilon(p)$ is well defined (for $k\ge2$ the
prefix has $\min<\max$, and $p_k$ is exactly one of
$\min-1$, $\max+1$). Conversely, given
$\varepsilon\in\{\sL,\sH\}^{n-1}$, set
$p_1=1+\#\{j:\varepsilon_j=\sL\}$ and iteratively
$p_k=(\text{current min})-1$ or $(\text{current max})+1$ according to
$\varepsilon_k$. This is the unique preimage and lands in $S_n$: the final
prefix is an interval of length $n$ equal to
$[\,p_1-\#\sL,\;p_1+\#\sH\,]=[1,n]$. The two maps are mutually inverse by
induction on $k$.

\eqref{eq:star}: if $\varepsilon_j=\sH$ then
$p_j=\max\{p_1,\dots,p_j\}>p_i$ for every $i<j$; if $\varepsilon_j=\sL$ then
$p_j<p_i$. The two cases are exhaustive and exclusive.
\end{proof}

\begin{remark}
The equivalence of (1)--(3) and the count $|\Av_n(132,312)|=2^{n-1}$ are
classical \cite{SimionSchmidt}; we have included the short proof to keep
the paper self-contained and because the reconstruction map and
rule~\eqref{eq:star} are used constantly below.
\end{remark}

\section{The characterization}\label{sec:thmA}

\begin{theorem}\label{thm:A}
Let $w=(s,p)$ be a nonnesting word with $p\in\Av_n(132,312)$ and sign word
$\varepsilon$. Then $w$ avoids $\{1132,3312\}$ if and only if
\[
 (C)\colon\qquad \varepsilon_K\neq\varepsilon_J
 \quad\text{for every pair } K<J \text{ with } q_1<o_J<q_K.
\]
\end{theorem}

Note that the pairs constrained by $(C)$ are crossing pairs of arcs
($K<J$ and $o_J<q_K$) whose later arc opens after the \emph{first} closer of
the word; automatically $K\ge2$, since $K=1$ would require $o_J<q_1$.

\begin{proof}
By Corollary~\ref{cor:combined}, $w$ contains a forbidden pattern iff there
is a triple $(v,x,y)$ with $v=p_i$ for some arc $i$ (every value is an arc
label, and $\snd(p_i)=q_i$ by Lemma~\ref{lem:fifo}), positions
$q_i<x<y$, and $w_y$ strictly between $p_i$ and $w_x$. By
Remark~\ref{rem:support}, $x$ and $y$ are letters of arcs $J,K>i$
respectively; and $J\neq K$, since strict betweenness forces
$w_x\neq w_y$.

\emph{Case $K>J$.} Then $i,J<K$, so $p_i,p_J\in\{p_1,\dots,p_K\}$, and by
Lemma~\ref{lem:signword}(2) the letter $w_y=p_K$ is the minimum or maximum
of that set --- never strictly between $p_i$ and $p_J$. So no violating
triple has $K>J$.

\emph{Case $K<J$ (hence $i<K<J$).} Arc $K$'s positions satisfy $o_K<o_J$ and
$q_K<q_J$. Since $x$ is a letter of $J$ and $y>x$ is a letter of $K$: if
$x=q_J$, then $y>q_J$, but both letters of arc $K$ sit at positions
$<q_J$ --- impossible; so $x=o_J$. If $y=o_K$, then $y=o_K<o_J=x$ ---
impossible; so $y=q_K$. The requirement $x<y$ reads $o_J<q_K$, and $x>q_i$
reads $q_i<o_J$. As for betweenness, rule~\eqref{eq:star} gives: if
$\varepsilon_J=\sH$ then $p_J>p_i$ and $p_J>p_K$, so $p_K$ is strictly
between $p_i$ and $p_J$ iff $p_i<p_K$ iff $\varepsilon_K=\sH$; if
$\varepsilon_J=\sL$ then $p_J$ lies below both, and betweenness holds iff
$p_K<p_i$ iff $\varepsilon_K=\sL$. In both cases:
\[
  p_K \text{ strictly between } p_i \text{ and } p_J
  \iff \varepsilon_K=\varepsilon_J .
\]
Crucially, this criterion does not depend on the witness arc $i$:
rule~\eqref{eq:star} gives $p_i<p_K\iff\varepsilon_K=\sH$ \emph{for every}
$i<K$.

Summarizing, a violating triple exists iff there are arcs $i<K<J$ with
$q_i<o_J<q_K$ and $\varepsilon_K=\varepsilon_J$. Finally, fix $K<J$ with
$o_J<q_K$; we claim an arc $i<K$ with $q_i<o_J$ exists iff $q_1<o_J$.
Indeed, if $q_1<o_J$ then $i=1$ works ($1<K$ holds because
$q_1<o_J<q_K$ forces $K\neq1$); conversely any such $i$ has
$q_1\le q_i<o_J$. Hence $w$ contains a forbidden pattern iff some pair
$K<J$ with $q_1<o_J<q_K$ has $\varepsilon_K=\varepsilon_J$, i.e.\ iff $(C)$
fails.
\end{proof}

\begin{corollary}\label{cor:pairs}
The map $w\mapsto(s(w),\varepsilon(p(w)))$ is a bijection from the set of
avoiders of $\nn$ onto the set of pairs $(s,\varepsilon)$ where $s$ is a
Dyck word of semilength $n$ and $\varepsilon\in\{\sL,\sH\}^{n-1}$ satisfies
$(C)$ for $s$. We call such $\varepsilon$ \emph{valid for $s$}.
\end{corollary}

\begin{proof}
Combine Lemma~\ref{lem:fifo} ($w\leftrightarrow(s,p)$, a bijection),
Lemma~\ref{lem:labels} (avoiders have $p\in\Av(132,312)$, so none are lost
in passing to sign words), Lemma~\ref{lem:signword}
($p\leftrightarrow\varepsilon$, a bijection), and Theorem~\ref{thm:A}
(avoidance $\iff(C)$). For $n=1$, $\varepsilon$ is the empty word, $(C)$ is
vacuous, and the unique word $11$ is an avoider.
\end{proof}

\begin{example}\label{ex:n3}
Let $n=3$ and $s=\sU\sU\sD\sU\sD\sD$, so $(o_1,o_2,o_3)=(1,2,4)$ and
$(q_1,q_2,q_3)=(3,5,6)$. The only pair $(K,J)$ with $q_1<o_J<q_K$ is
$(2,3)$: $q_1=3<o_3=4<q_2=5$. So $(C)$ reads
$\varepsilon_2\neq\varepsilon_3$, with the two solutions
$\varepsilon\in\{\sL\sH,\sH\sL\}$, i.e.\ $p\in\{(2,1,3),(2,3,1)\}$. The word
pattern of this shape is $p_1p_2p_1p_3p_2p_3$, so the two avoiders are
$212313$ and $232131$. Indeed the other four label permutations all fail;
e.g.\ $p=(1,2,3)$ gives $121323$, which contains $1132$ as the subsequence
$1\,1\,3\,2$ at positions $1,3,4,5$. This matches Theorem~\ref{thm:A}, and
the count $2$ will reappear in Example~\ref{ex:table} as the two codes
$\sC\sC\sA,\sC\sC\sB$.
\end{example}

\section{The admissible shapes}\label{sec:shapes}

Fix a Dyck word $s$ of semilength $n$. Let $a=a(s)\ge1$ be the length of its
first ascent run, so that $o_j=j$ for $j\le a$ and $q_1=a+1$. Call an arc
$J$ \emph{late} if $o_J>q_1$. For a late arc $J$ let
\[
  S_J \;=\; \{K<J : o_J<q_K\},
\]
so that the pairs constrained by $(C)$ are exactly the pairs $(K,J)$ with
$J$ late and $K\in S_J$. (For $K<J$ we have $o_K<o_J$ automatically, so
$S_J$ is the set of arcs \emph{open} at the moment $o_J$, i.e.\ with
$o_K<o_J<q_K$.)

\begin{lemma}\label{lem:interval}
Let $J\in[n]$ and let $c$ be the number of closers at positions $<o_J$.
Then $S_J=\{c+1,\dots,J-1\}$, and $|S_J|=J-1-c$ equals the height of $s$
(number of $\sU$'s minus number of $\sD$'s) just before position $o_J$.
Moreover, $J$ is late if and only if $J\ge a+1$.
\end{lemma}

\begin{proof}
The arcs opened before $o_J$ are exactly $1,\dots,J-1$ (openers are
sorted). The closers before $o_J$ are an initial segment of
$q_1<q_2<\cdots$, hence are exactly $q_1,\dots,q_c$: arcs $1,\dots,c$ are
closed before $o_J$ and arcs $c+1,\dots,J-1$ remain open ($c\le J-1$ since
each closed arc was opened). The height before a position equals the number
of earlier openers minus earlier closers, i.e.\ $(J-1)-c$. Finally, openers
$1,\dots,a$ sit at positions $1,\dots,a<q_1=a+1$, and every opener $o_J$ with
$J\ge a+1$ satisfies $o_J>a+1=q_1$ because position $a+1$ is a closer.
\end{proof}

Define, with respect to the decomposition above:
\begin{itemize}
\item a \emph{gap opener}: an opener at a position strictly between $q_1$
  and $q_a$ (for $a=1$ there is none, as the interval is empty);
\item a \emph{tail opener}: an opener at a position $>q_a$; the
  corresponding arcs are \emph{tail arcs}. For a tail arc $j$, let
  $\delta_j\in\{0,1,2,\dots\}$ be the height of $s$ just before position
  $o_j$ (its \emph{start height}).
\end{itemize}

\begin{lemma}[obstructions]\label{lem:obstructions}
If $s$ has at least two gap openers, or a tail opener with start height
$\ge2$, then no $\varepsilon\in\{\sL,\sH\}^{n-1}$ is valid for $s$.
\end{lemma}

\begin{proof}
\emph{Two gap openers.} Then $a\ge2$. The openers in the interval
$(q_1,q_a)$ belong to arcs $a+1,a+2,\dots$ in order (openers are sorted and
arcs $\le a$ open before $q_1$); let $\beta=a+1$ and $\beta'=a+2$ be the
arcs of the first two. Both are late. Let $c\le c'$ be the numbers of
closers before $o_\beta$ and $o_{\beta'}$; since
$q_1<o_\beta<o_{\beta'}<q_a$ we get $1\le c\le c'\le a-1$. By
Lemma~\ref{lem:interval}, $S_\beta=\{c+1,\dots,a\}\ni a$ and
$S_{\beta'}=\{c'+1,\dots,a+1\}\supseteq\{a,\beta\}$. So the pairs
$(a,\beta)$, $(a,\beta')$, $(\beta,\beta')$ are all constrained by $(C)$,
and a valid $\varepsilon$ would need
$\varepsilon_a,\varepsilon_\beta,\varepsilon_{\beta'}$ pairwise distinct in
the two-element set $\{\sL,\sH\}$ --- impossible.

\emph{Tail opener at height $\ge2$.} Let $J$ be its arc, so $o_J>q_a$ and,
by Lemma~\ref{lem:interval}, $S_J=\{c+1,\dots,J-1\}$ with $J-1-c\ge2$;
hence $K:=J-2$ and $K':=J-1$ both lie in $S_J$. Note $K\ge c+1\ge a+1\ge2$,
since at least the $a$ closers $q_1,\dots,q_a$ precede $o_J$. The arc $K'$
is late: $q_{K'}>o_J>q_a$ forces $K'>a$ (closers are sorted), so
$K'\ge a+1$. Also $K\in S_{K'}$: $K<K'$ and $q_K>o_J>o_{K'}$. So the pairs
$(K,K')$, $(K,J)$, $(K',J)$ are all constrained by $(C)$, again requiring
three pairwise distinct values in $\{\sL,\sH\}$ --- impossible.
\end{proof}

It remains to describe the shapes with at most one gap opener and all tail
start heights $\le1$; we call these \emph{admissible}. For $m\ge0$,
$h_0\in\{0,1\}$ and $d_1,\dots,d_m\in\{0,1\}$ with $d_1\le h_0$ when
$m\ge1$, let $\mathrm{Tail}(h_0;\,d_1,\dots,d_m)$ be the $\{\sU,\sD\}$-word
produced by the following scan: start with $h:=h_0$; for $r=1,\dots,m$ emit
$\sD^{\,h-d_r}$ then $\sU$, and set $h:=d_r+1$; finally emit $\sD^{\,h}$.
(All $\sD$-run lengths are $\ge0$: $d_1\le h_0$, and for $r\ge2$,
$d_r\le1\le d_{r-1}+1=h$. For $m=0$ the word is $\sD^{\,h_0}$.) By
construction, $\mathrm{Tail}$ describes a lattice path from height $h_0$ to
height $0$, never negative, whose $r$th $\sU$-step starts at height $d_r$.

\begin{lemma}[classification]\label{lem:classification}
Let $s$ be admissible. Then exactly one of the following holds, in each case
with the stated parameters determined by $s$; and conversely, every choice
of parameters yields an admissible Dyck word of semilength $n$ with those
data.
\begin{itemize}
\item[(S)] \emph{Staircase}: $a=n$ and $s=\sU^n\sD^n$. (No parameters.)
\item[(I)] $1\le a\le n-1$ and $s$ has no gap opener. The tail arcs are
  $a+1,\dots,n$; $\delta_{a+1}=0$; and
  \[
    s=\sU^a\sD^a\cdot\mathrm{Tail}\bigl(0;\,\delta_{a+1},\dots,\delta_n\bigr),
  \]
  with parameters $\bigl(a,\ \delta=(\delta_{a+2},\dots,\delta_n)\in\{0,1\}^{n-a-1}\bigr)$
  and $\delta_{a+1}:=0$.
\item[(II)] $2\le a\le n-1$ and $s$ has exactly one gap opener, which
  belongs to the arc $B:=a+1$ and has exactly $i$ closers before it, where
  $1\le i\le a-1$; moreover $S_B=\{i+1,\dots,a\}$. The tail arcs are
  $a+2,\dots,n$, and
  \[
    s=\sU^a\sD^{\,i}\,\sU\,\sD^{\,a-i}\cdot
      \mathrm{Tail}\bigl(1;\,\delta_{a+2},\dots,\delta_n\bigr),
  \]
  with parameters $\bigl(a,\,i,\ \delta=(\delta_{a+2},\dots,\delta_n)\in\{0,1\}^{n-a-1}\bigr)$.
\end{itemize}
\end{lemma}

\begin{proof}
If $a=n$ then all $n$ arcs open in the first run and $s=\sU^n\sD^n$: case
(S). Assume $a<n$.

\emph{No gap opener.} Every position strictly between $q_1$ and $q_a$ is a
closer, so $q_m=a+m$ for $m\le a$ and the word begins $\sU^a\sD^a$. All
arcs $\ge a+1$ open after $q_a=2a$: the tail arcs are exactly
$a+1,\dots,n$, a nonempty set since $a<n$. The height after position $2a$
is $0$ and the next position must be an opener (a closer would make the
height negative), so $\delta_{a+1}=0$; by admissibility $\delta_j\le1$ for
the rest. The remainder of the word is forced by the start heights:
between consecutive tail openers all positions are closers, and the height
must drop from $\delta_j+1$ (just after the $\sU$ of tail arc $j$) to
$\delta_{j+1}$ (just before the $\sU$ of tail arc $j+1$), giving exactly
$\delta_j+1-\delta_{j+1}$ letters $\sD$; after the last opener the height
$\delta_n+1$ returns to $0$. This is precisely
$\mathrm{Tail}(0;\delta_{a+1},\dots,\delta_n)$, so $s$ has the displayed
form and $(a,\delta)$ is determined by $s$. Conversely, for any
$(a,\delta)$ the displayed word is a Dyck word of semilength $n$
($\sU$-count $a+(n-a)=n$; the scan keeps $h\ge0$ and ends at height $0$),
its first ascent is exactly $a$ (position $a+1$ is a $\sD$ since $a\ge1$),
it has no gap opener (positions $a+1,\dots,2a$ are $\sD$'s, and $q_a=2a$),
and its tail start heights are $\delta_{a+1},\dots,\delta_n$ by the
construction of the scan: case (I) with the same parameters.

\emph{One gap opener.} The gap opener is the first opener after $q_1$,
hence belongs to arc $a+1=:B$; the interval $(q_1,q_a)$ must be nonempty,
so $a\ge2$. Let $i$ be the number of closers before $o_B$: $o_B>q_1$ gives
$i\ge1$ and $o_B<q_a$ gives $i\le a-1$. Positions in $(q_1,q_a)$ other than
$o_B$ are closers, so the prefix of $s$ up to $q_a$ reads
$\sU^a\sD^{\,i}\sU\sD^{\,a-i}$ (closers $q_1,\dots,q_i$, then $o_B$, then
$q_{i+1},\dots,q_a$); in particular $q_a=2a+1$ and, by
Lemma~\ref{lem:interval} with $c=i$, $S_B=\{i+1,\dots,a\}$. The height
after $q_a$ is $(a+1)-a=1$. The tail arcs are $a+2,\dots,n$ (possibly
none), and the same forced-run argument as in case (I) identifies the
remainder with $\mathrm{Tail}(1;\delta_{a+2},\dots,\delta_n)$; for $a=n-1$
this is $\mathrm{Tail}(1;\,)=\sD$. So $(a,i,\delta)$ is determined by $s$,
and conversely each $(a,i,\delta)$ yields a Dyck word of case (II) by the
same checks ($\sU$-count $a+1+(n-a-1)=n$; first ascent $a$ since $i\ge1$
puts a $\sD$ at position $a+1$; exactly one gap opener, at position
$a+1+i<2a+1=q_a$, with $i$ closers before it; tail start heights
$\delta_j$).

The three cases are pairwise exclusive ($a=n$ versus $a<n$; gap opener
absent versus present), and every admissible shape falls in one of them.
\end{proof}

\begin{lemma}[local form of validity]\label{lem:inventory}
Let $s$ be admissible and $\varepsilon\in\{\sL,\sH\}^{n-1}$ (indexed
$2,\dots,n$). Then $\varepsilon$ is valid for $s$ if and only if:
\begin{itemize}
\item[(S)] always (there are no late arcs, hence no constrained pairs);
\item[(I)] for every tail arc $j\ge a+2$ with $\delta_j=1$:
  $\varepsilon_j=\overline{\varepsilon_{j-1}}$;
\item[(II)] $\varepsilon_{i+1}=\cdots=\varepsilon_a
  =\overline{\varepsilon_{a+1}}$, and for every tail arc $j\ge a+2$ with
  $\delta_j=1$: $\varepsilon_j=\overline{\varepsilon_{j-1}}$.
\end{itemize}
\end{lemma}

\begin{proof}
Validity is the conjunction, over late arcs $J$ (i.e.\ $J\ge a+1$, by
Lemma~\ref{lem:interval}) and $K\in S_J$, of
$\varepsilon_K\neq\varepsilon_J$. We list the sets $S_J$ case by case
using Lemma~\ref{lem:interval}.

(S): $a=n$ and there are no late arcs; the conjunction is empty.

(I): the late arcs are $a+1,\dots,n$, all tail arcs. For $J=a+1$: the
closers before $o_{a+1}=2a+1$ are $q_1,\dots,q_a$, so
$S_{a+1}=\{a+1,\dots,a\}=\varnothing$ --- no constraint. For a tail arc
$J\ge a+2$: $|S_J|=\delta_J\in\{0,1\}$ by Lemma~\ref{lem:interval}. If
$\delta_J=0$ there is no constraint; if $\delta_J=1$ then $S_J$, being the
interval $\{c+1,\dots,J-1\}$ of size one, equals $\{J-1\}$, and the
constraint $\varepsilon_J\neq\varepsilon_{J-1}$ over a two-letter alphabet
reads $\varepsilon_J=\overline{\varepsilon_{J-1}}$.

(II): the late arcs are $a+1,\dots,n$. For $J=B=a+1$:
$S_B=\{i+1,\dots,a\}$ by Lemma~\ref{lem:classification}, and the
constraints $\{\varepsilon_K\neq\varepsilon_B: K=i+1,\dots,a\}$ hold iff
$\varepsilon_{i+1}=\cdots=\varepsilon_a=\overline{\varepsilon_{a+1}}$ (the
block is nonempty as $i\le a-1$). For tail arcs $J\ge a+2$ the analysis is
exactly as in (I); note that for $J=a+2$ with $\delta_J=1$ the constraint
set is $S_J=\{a+1\}=\{B\}$, also covered by the rule
$\varepsilon_J=\overline{\varepsilon_{J-1}}$.
\end{proof}

\begin{corollary}[free coordinates]\label{cor:free}
For each admissible $s$ define the \emph{free set} $F(s)\subseteq\{2,\dots,n\}$:
\[
  F(s)=
  \begin{cases}
   \{2..n\} & \text{(S)},\\
   \{2..a+1\}\cup\{j\ge a+2:\delta_j=0\} & \text{(I)},\\
   \{2..i\}\cup\{a+1\}\cup\{j\ge a+2:\delta_j=0\} & \text{(II)}.
  \end{cases}
\]
Then the restriction map $\varepsilon\mapsto\varepsilon|_{F(s)}$ is a
bijection from the valid sign words for $s$ onto $\{\sL,\sH\}^{F(s)}$. In
particular the number of valid sign words is $N(s)=2^{|F(s)|}$, with
\[
  |F(s)| =
  \begin{cases}
   n-1 & \text{(S)},\\
   a+z & \text{(I)},\\
   i+z & \text{(II)},
  \end{cases}
  \qquad\text{where } z:=\#\{a+2\le j\le n:\delta_j=0\}.
\]
\end{corollary}

\begin{proof}
Define an extension map $E$ on $f\in\{\sL,\sH\}^{F(s)}$: set
$\varepsilon_j:=f_j$ for $j\in F(s)$; in case (II) set
$\varepsilon_{i+1}=\cdots=\varepsilon_a:=\overline{f_{a+1}}$; then process
the tail arcs $j=a+2,\dots,n$ in increasing order, setting
$\varepsilon_j:=\overline{\varepsilon_{j-1}}$ whenever $\delta_j=1$. Each
reference $\varepsilon_{j-1}$ is already defined at the time it is used:
$j-1$ is in $F(s)$, or in the block $\{i+1..a\}$, or equals $a+1$, or is an
earlier-processed tail arc. By construction $E(f)$ satisfies the conditions
of Lemma~\ref{lem:inventory}, hence is valid, and $E(f)|_{F(s)}=f$.
Conversely, if $\varepsilon$ is valid, then by Lemma~\ref{lem:inventory}
its non-free coordinates are determined from $\varepsilon|_{F(s)}$ by
exactly the assignments that $E$ makes --- the block coordinates equal
$\overline{\varepsilon_{a+1}}$, and each $\delta_j=1$ tail coordinate
equals $\overline{\varepsilon_{j-1}}$, by induction on $j$ --- so
$E(\varepsilon|_{F(s)})=\varepsilon$. Thus restriction and extension are
mutually inverse. The values of $|F(s)|$ are immediate counts: in case (S),
$|\{2..n\}|=n-1$; in case (I), $|\{2..a+1\}|=a$, giving $a+z$; in case
(II), $|\{2..i\}|=i-1$ plus the singleton $\{a+1\}$, giving $i+z$.
\end{proof}

\section{First proof of the main theorem: summation}\label{sec:count}

\begin{theorem}\label{thm:main}
For every $n\ge1$, the number of nonnesting permutations of $\nn$ avoiding
$1132$ and $3312$ is
\[
  c_n \;=\; 3^n-3\cdot2^{n-1}+1 .
\]
\end{theorem}

\begin{proof}
By Corollary~\ref{cor:pairs}, $c_n=\sum_s N(s)$, the sum over all Dyck
words $s$ of semilength $n$ of the number of valid sign words. By
Lemma~\ref{lem:obstructions}, $N(s)=0$ unless $s$ is admissible, and by
Lemma~\ref{lem:classification} the admissible shapes are parametrized
bijectively by the data (S), $(a,\delta)$ in case (I) with $1\le a\le n-1$,
and $(a,i,\delta)$ in case (II) with $2\le a\le n-1$, $1\le i\le a-1$,
where $\delta\in\{0,1\}^{n-a-1}$ in both cases. By
Corollary~\ref{cor:free}, $N(s)=2^{|F(s)|}$ with the listed exponents. We
sum case by case. Since each coordinate $\delta_j$ ranges freely over
$\{0,1\}$ and contributes a factor $2$ to $2^{|F(s)|}$ when $\delta_j=0$
and $1$ when $\delta_j=1$,
\[
 \sum_{\delta\in\{0,1\}^{n-a-1}} 2^{\,z(\delta)} \;=\; (2+1)^{\,n-a-1} \;=\;
 3^{\,n-a-1}.
\]
Hence
\[
  c_n \;=\; \underbrace{2^{n-1}}_{\text{(S)}}
  \;+\; \underbrace{\sum_{a=1}^{n-1} 2^{a}\,3^{\,n-1-a}}_{\text{(I)}}
  \;+\; \underbrace{\sum_{a=2}^{n-1}\Bigl(\sum_{i=1}^{a-1}2^{i}\Bigr)3^{\,n-1-a}}_{\text{(II)}}
  \;=\; 2^{n-1} \;+\; \Sigma_1 \;+\; \Sigma_2,
\]
with $\sum_{i=1}^{a-1}2^i=2^a-2$. We use the identity
\begin{equation}\label{eq:geom}
  \sum_{a=1}^{m} 2^{a-1}3^{\,m-a} \;=\; 3^m-2^m \qquad (m\ge0),
\end{equation}
proved by induction on $m$: both sides vanish at $m=0$ and satisfy
$g(m)=3g(m-1)+2^{m-1}$. Applying \eqref{eq:geom} with $m=n-1$,
\[
  \Sigma_1=2\sum_{a=1}^{n-1}2^{a-1}3^{(n-1)-a}=2\bigl(3^{n-1}-2^{n-1}\bigr).
\]
For $\Sigma_2$, assume first that $n\ge3$. Removing the $a=1$ term
$2\cdot3^{n-2}$ from $\Sigma_1$ and using
$\sum_{a=2}^{n-1}3^{n-1-a}=\sum_{k=0}^{n-3}3^k=\tfrac{3^{n-2}-1}{2}$,
\[
  \Sigma_2=\Bigl(2\bigl(3^{n-1}-2^{n-1}\bigr)-2\cdot3^{n-2}\Bigr)
  -2\cdot\frac{3^{n-2}-1}{2}
  =2\cdot3^{n-1}-3\cdot 3^{n-2}-2^{n}+1
  =3^{n-1}-2^{n}+1 .
\]
For $n\in\{1,2\}$ the sum $\Sigma_2$ is empty and the expression
$3^{n-1}-2^n+1$ equals $0$, so this closed form holds for all $n\ge1$.
Adding up,
\[
  c_n=2^{n-1}+2\cdot3^{n-1}-2^{n}+3^{n-1}-2^{n}+1
     =3^{n}+2^{n-1}-2^{n+1}+1
     =3^{n}-3\cdot2^{n-1}+1 . \qedhere
\]
\end{proof}

\begin{remark}\label{rem:structure}
The three exponential scales of the answer are transparent in the
decomposition: each free tail coordinate contributes a factor $3$ (the
choices ``$\delta_j=0$ with $\varepsilon_j=\sL$'', ``$\delta_j=0$ with
$\varepsilon_j=\sH$'', ``$\delta_j=1$, sign forced''), the first-run
coordinates contribute factors of $2$, and the bookkeeping of the three
families produces exactly $3^n-3\cdot2^{n-1}+1$. This matches the ordinary
generating function
$x(1-2x+3x^2)/\bigl((1-x)(1-2x)(1-3x)\bigr)$
of the (shifted) OEIS entry A168583 \cite{OEIS}, with poles at
$1,\tfrac12,\tfrac13$.
\end{remark}

\section{Second proof of the endgame: an explicit bijection}\label{sec:bijection}

The summation in Section~\ref{sec:count} can be replaced by an explicit
bijection: the classification data $(a,i,\delta,\varepsilon|_{F(s)})$ of an
avoider can be packaged into a single ternary string, in such a way that
the image language has a transparent description and count. Throughout this
section, fix the alphabet $\{\sA,\sB,\sC\}$ and identify signs with letters
via
\[
 \ell(\sL)=\sA,\qquad \ell(\sH)=\sB .
\]

\begin{definition}\label{def:Wn}
$\Wn\subseteq\{\sA,\sB,\sC\}^n$ is the set of strings $\sigma$ such that:
\begin{itemize}
\item if $\sigma$ contains no $\sC$, then $\sigma_1=\sA$;
\item if $\sigma_1=\sC$, then there are positions $2\le r<t\le n$ with
  $\sigma_r=\sC$ and $\sigma_t\neq\sC$;
\item (if $\sigma_1\neq\sC$ and $\sigma$ contains a $\sC$: no condition).
\end{itemize}
\end{definition}

\begin{lemma}\label{lem:Wcount}
The complement of $\Wn$ in $\{\sA,\sB,\sC\}^n$ is the disjoint union of
\[
 X_1=\sB\,\{\sA,\sB\}^{n-1}
 \quad\text{and}\quad
 X_2=\{\,\sC\,v\,\sC^{\,k}\;:\; v\in\{\sA,\sB\}^{m},\ m,k\ge0,\
 1+m+k=n \,\},
\]
with $|X_1|=2^{n-1}$ and $|X_2|=\sum_{m=0}^{n-1}2^m=2^n-1$. Hence
\[
  |\Wn| \;=\; 3^n-2^{n-1}-(2^n-1) \;=\; 3^n-3\cdot2^{n-1}+1
  \qquad(n\ge1).
\]
\end{lemma}

\begin{proof}
A string violates Definition~\ref{def:Wn} iff either it is $\sC$-free with
$\sigma_1=\sB$ (a $\sC$-free string cannot have $\sigma_1=\sC$), or
$\sigma_1=\sC$ and no $\sC$ at a position $r\ge2$ is followed later by a
non-$\sC$. The first family is $X_1$. For the second, write
$\sigma=\sC\rho$: the condition says that every letter after the first
$\sC$ of $\rho$ (if any) is again $\sC$, i.e.\
$\rho\in\{\sA,\sB\}^*\sC^*$, which is exactly the family $X_2$ (the case
$k=0$ covering $\sC$-free $\rho$). $X_1$ and $X_2$ are disjoint (first
letter $\sB$ versus $\sC$), and within $X_2$ the pair $(m,k)$ is determined
by $\sigma$ ($v$ is the maximal $\{\sA,\sB\}$-prefix of $\rho$), so
$|X_2|=\sum_{m+k=n-1}2^m=2^n-1$.
\end{proof}

\begin{definition}[the encoding $\Phi$]\label{def:phi}
Let $s$ be admissible and $\varepsilon$ valid for $s$. Define
$\sigma=\Phi(s,\varepsilon)\in\{\sA,\sB,\sC\}^n$ by cases on
Lemma~\ref{lem:classification}:
\begin{itemize}
\item[(S)] $\sigma_1=\sA$; and $\sigma_j=\ell(\varepsilon_j)$ for
 $2\le j\le n$.
\item[(I)] $\sigma_j=\ell(\varepsilon_{j+1})$ for $1\le j\le a$;\quad
 $\sigma_{a+1}=\sC$;\quad and for $a+2\le j\le n$:
 $\sigma_j=\sC$ if $\delta_j=1$, else $\sigma_j=\ell(\varepsilon_j)$.
\item[(II)] $\sigma_1=\sC$;\quad $\sigma_j=\ell(\varepsilon_j)$ for
 $2\le j\le i$;\quad $\sigma_j=\sC$ for $i+1\le j\le a$;\quad
 $\sigma_{a+1}=\ell(\varepsilon_{a+1})$;\quad and for $a+2\le j\le n$:
 $\sigma_j=\sC$ if $\delta_j=1$, else $\sigma_j=\ell(\varepsilon_j)$.
\end{itemize}
\end{definition}

In words: every free coordinate of $\varepsilon$ is written as its sign
letter, every \emph{forced} arc (the block $\{i+1..a\}$ of case (II), and
the $\delta_j=1$ tail arcs) is written $\sC$; in case (I) the initial free
block $\varepsilon_2,\dots,\varepsilon_{a+1}$ is shifted one slot to the
left so that position $a+1$ can carry a $\sC$ marking the end of the first
descent; in case (II), position $1$ carries a $\sC$ marking the presence of
the gap opener; in case (S), position $1$ carries the padding letter
$\sA$.

\begin{theorem}\label{thm:B}
$\Phi$ is a bijection from
$\{(s,\varepsilon): s \text{ admissible},\ \varepsilon \text{ valid for } s\}$
onto $\Wn$.
\end{theorem}

\begin{proof}
\textbf{Step 1: $\Phi$ lands in $\Wn$, with disjoint images for the three
cases.} (S): $\sigma$ is $\sC$-free with $\sigma_1=\sA$, so $\sigma\in\Wn$.
(I): $\sigma_1=\ell(\varepsilon_2)\in\{\sA,\sB\}$ and $\sigma$ contains a
$\sC$ (at position $a+1\le n$), so $\sigma\in\Wn$ by the third clause.
(II): $\sigma_1=\sC$; take $r=i+1$ and $t=a+1$: then $2\le r<t\le n$,
$\sigma_r=\sC$ (the block is nonempty since $i\le a-1$) and
$\sigma_t\in\{\sA,\sB\}$, so $\sigma\in\Wn$ by the second clause.
Disjointness: (S)-images are $\sC$-free; (I)-images contain a $\sC$ and
have $\sigma_1\neq\sC$; (II)-images have $\sigma_1=\sC$.

\textbf{Step 2: a parse map $\Psi$ defined on all of $\Wn$.} Let
$\sigma\in\Wn$.

\emph{Case $\sC\notin\sigma$.} Then $\sigma_1=\sA$ by
Definition~\ref{def:Wn}. Output case (S) with
$\varepsilon_j:=\ell^{-1}(\sigma_j)$ for $2\le j\le n$.

\emph{Case $\sigma_1\neq\sC$, $\sC\in\sigma$.} Let $a+1$ be the position of
the first $\sC$, so $1\le a\le n-1$ and
$\sigma_1,\dots,\sigma_a\in\{\sA,\sB\}$. Set
$\varepsilon_{j+1}:=\ell^{-1}(\sigma_j)$ for $1\le j\le a$; then for
$j=a+2,\dots,n$ in increasing order: if $\sigma_j=\sC$ set $\delta_j:=1$
and $\varepsilon_j:=\overline{\varepsilon_{j-1}}$, else set $\delta_j:=0$
and $\varepsilon_j:=\ell^{-1}(\sigma_j)$. Output case (I) with parameters
$(a,\delta)$ --- an admissible shape by the converse direction of
Lemma~\ref{lem:classification}(I) --- and the sign word $\varepsilon$,
which is valid by Lemma~\ref{lem:inventory}(I), since its $\delta_j=1$
coordinates satisfy the forced equalities by construction.

\emph{Case $\sigma_1=\sC$.} Let $i-1\ge0$ be the length of the maximal
$\{\sA,\sB\}$-run starting at position $2$, so positions $2,\dots,i$ carry
letters in $\{\sA,\sB\}$ and position $i+1$, if it exists, carries $\sC$.
The second clause of Definition~\ref{def:Wn} provides $2\le r<t\le n$ with
$\sigma_r=\sC$ and $\sigma_t\neq\sC$; in particular there is a $\sC$ at a
position $\ge2$, so $i+1\le n$ and $\sigma_{i+1}=\sC$. Let the maximal
$\sC$-run starting at position $i+1$ occupy positions $i+1,\dots,a$. We
claim $a\le n-1$ and $\sigma_{a+1}\in\{\sA,\sB\}$. Consider the witness
pair $(r,t)$, and note $r\ge i+1$ (positions $2,\dots,i$ are not $\sC$). If
$r\le a$: then $t>r$ and $\sigma_t\neq\sC$ force $t\ge a+1$ (positions
$r+1,\dots,a$ are $\sC$), so position $a+1$ exists, and
$\sigma_{a+1}\neq\sC$ by maximality of the run. If $r>a$: then
$a+1\le r\le n$, so position $a+1$ exists; $r=a+1$ is impossible (a $\sC$
at $a+1$ would extend the maximal run), so again $\sigma_{a+1}\neq\sC$ by
maximality. This proves the claim. Now set
$\varepsilon_j:=\ell^{-1}(\sigma_j)$ for $2\le j\le i$;
$\varepsilon_{a+1}:=\ell^{-1}(\sigma_{a+1})$;
$\varepsilon_{i+1}=\cdots=\varepsilon_a:=\overline{\varepsilon_{a+1}}$; and
process $j=a+2,\dots,n$ exactly as in the previous case (defining
$\delta_j$ and $\varepsilon_j$). Output case (II) with parameters
$(a,i,\delta)$ --- well defined, since $1\le i\le a-1\le n-2$, and an
admissible shape by Lemma~\ref{lem:classification}(II) --- and the sign
word $\varepsilon$, valid by Lemma~\ref{lem:inventory}(II).

\textbf{Step 3: $\Psi\circ\Phi=\mathrm{id}$.} Let $(s,\varepsilon)$ be a
valid pair and $\sigma=\Phi(s,\varepsilon)$.
\begin{itemize}
\item (S): $\sigma$ is $\sC$-free, so $\Psi$ takes the (S)-branch and
 recovers each $\varepsilon_j$ from $\sigma_j$.
\item (I): $\sigma_1\neq\sC$, and the first $\sC$ of $\sigma$ is at
 position $a+1$, since positions $1,\dots,a$ carry sign letters. So $\Psi$
 takes the (I)-branch and recovers $a$; it recovers
 $\varepsilon_2,\dots,\varepsilon_{a+1}$ from positions $1,\dots,a$; and
 for $j\ge a+2$ it recovers $\delta_j$ ($\sC\leftrightarrow\delta_j=1$) and
 $\varepsilon_j$ --- directly if $\delta_j=0$, and as
 $\overline{\varepsilon_{j-1}}$ if $\delta_j=1$, which equals the original
 $\varepsilon_j$ because the original is valid
 (Lemma~\ref{lem:inventory}(I)); the argument is an induction on $j$ using
 that $\varepsilon_{j-1}$ was recovered correctly. By
 Lemma~\ref{lem:classification}(I), the parameters $(a,\delta)$ determine
 $s$.
\item (II): $\sigma_1=\sC$, so $\Psi$ takes the (II)-branch. In $\sigma$,
 positions $2,\dots,i$ carry sign letters, positions $i+1,\dots,a$ carry
 $\sC$ (a nonempty block), and position $a+1$ carries a sign letter; hence
 the maximal $\{\sA,\sB\}$-run from position $2$ ends exactly at $i$ and
 the maximal $\sC$-run from $i+1$ ends exactly at $a$: $\Psi$ recovers
 $(i,a)$. It then recovers $\varepsilon_2,\dots,\varepsilon_i$ and
 $\varepsilon_{a+1}$ directly; the block coordinates as
 $\overline{\varepsilon_{a+1}}$ --- correct by
 Lemma~\ref{lem:inventory}(II) --- and the tail as in case (I). By
 Lemma~\ref{lem:classification}(II), $(a,i,\delta)$ determine $s$.
\end{itemize}

\textbf{Step 4: $\Phi\circ\Psi=\mathrm{id}$.} Let $\sigma\in\Wn$ and
$(s,\varepsilon)=\Psi(\sigma)$. In each branch, the output of $\Psi$
classifies into the same case it was parsed in (this is the converse
direction of Lemma~\ref{lem:classification}, which asserts that the
constructed shape lies in case (I), resp.\ (II), with the same parameters),
and $\Phi$ applied to it writes back exactly the letters read: in (S), the
letter $\sA$ followed by the sign letters; in (I), the $a$ sign letters
read from positions $1,\dots,a$, the $\sC$ at position $a+1$, and for
$j\ge a+2$ a $\sC$ iff $\delta_j=1$ iff $\sigma_j=\sC$, else
$\ell(\varepsilon_j)=\sigma_j$; in (II), the leading $\sC$, the run
structure (positions $2,\dots,i$; the block $i+1,\dots,a$; position
$a+1$), and the tail, likewise. So $\Phi(\Psi(\sigma))=\sigma$.

By Steps 1--4, $\Phi$ is a bijection with inverse $\Psi$.
\end{proof}

\begin{corollary}[second proof of Theorem~\ref{thm:main}]\label{cor:secondproof}
Composing the bijections
\[
 \{\text{avoiders of }\nn\}
 \xrightarrow{\ \text{Cor.~\ref{cor:pairs}}\ }
 \{(s,\varepsilon):\varepsilon\text{ valid for }s\}
 \xrightarrow{\ \Phi\ } \Wn
\]
(the middle set ranges over all Dyck words, but by
Lemmas~\ref{lem:obstructions} and~\ref{lem:classification} only admissible
shapes carry valid sign words, so the middle set is exactly the domain of
$\Phi$) and applying Lemma~\ref{lem:Wcount} gives
$c_n=|\Wn|=3^n-3\cdot2^{n-1}+1$. \qed
\end{corollary}

The end-to-end map is explicit and runs in linear time in both directions:
given an avoider $w$, read off $(s,p)$ (Lemma~\ref{lem:fifo}), convert $p$
to $\varepsilon$ (Lemma~\ref{lem:signword}), and write $\sigma$
(Definition~\ref{def:phi}); given $\sigma\in\Wn$, parse it
(Theorem~\ref{thm:B}, Step 2), rebuild $s$ from the parameters
(Lemma~\ref{lem:classification}), rebuild $p$ from $\varepsilon$
(Lemma~\ref{lem:signword}), and interleave (Lemma~\ref{lem:fifo}).

\begin{example}\label{ex:table}
For $n=3$, the $16$ avoiders and their codes
$\sigma=\Phi(\,\cdot\,)$ are:
\begin{center}
\begin{tabular}{cll@{\qquad}cll}
\toprule
$\sigma$ & word & case & $\sigma$ & word & case\\
\midrule
$\sA\sA\sA$ & $321321$ & S & $\sA\sC\sB$ & $221133$ & I, $a=1$\\
$\sA\sA\sB$ & $213213$ & S & $\sA\sC\sC$ & $221313$ & I, $a=1$\\
$\sA\sA\sC$ & $323211$ & I, $a=2$ & $\sB\sA\sC$ & $232311$ & I, $a=2$\\
$\sA\sB\sA$ & $231231$ & S & $\sB\sB\sC$ & $121233$ & I, $a=2$\\
$\sA\sB\sB$ & $123123$ & S & $\sB\sC\sA$ & $223311$ & I, $a=1$\\
$\sA\sB\sC$ & $212133$ & I, $a=2$ & $\sB\sC\sB$ & $112233$ & I, $a=1$\\
$\sA\sC\sA$ & $332211$ & I, $a=1$ & $\sB\sC\sC$ & $223131$ & I, $a=1$\\
 & & & $\sC\sC\sA$ & $232131$ & II\\
 & & & $\sC\sC\sB$ & $212313$ & II\\
\bottomrule
\end{tabular}
\end{center}
The $3\cdot2^{2}-1=11$ excluded strings are
$X_1=\{\sB\sA\sA,\sB\sA\sB,\sB\sB\sA,\sB\sB\sB\}$ and
$X_2=\{\sC\sA\sA$, $\sC\sA\sB$, $\sC\sB\sA$, $\sC\sB\sB$, $\sC\sA\sC$,
$\sC\sB\sC$, $\sC\sC\sC\}$.
\end{example}

\begin{example}\label{ex:cbccac}
A case (II) example with $n=6$: $\sigma=\sC\sB\sC\sC\sA\sC$. Parsing:
$\sigma_1=\sC$; the $\{\sA,\sB\}$-run from position $2$ is ``$\sB$'', so
$i=2$ and $\varepsilon_2=\sH$; the $\sC$-run occupies positions $3,4$, so
$a=4$; $\varepsilon_5=\ell^{-1}(\sA)=\sL$; the block gives
$\varepsilon_3=\varepsilon_4=\overline{\sL}=\sH$; the tail arc $6$ has
$\sigma_6=\sC$, so $\delta_6=1$ and
$\varepsilon_6=\overline{\varepsilon_5}=\sH$. The shape is
$s=\sU^4\sD^2\sU\sD^2\cdot\mathrm{Tail}(1;\,1)
 =\sU\sU\sU\sU\sD\sD\sU\sD\sD\sU\sD\sD$, the sign word
$(\sH,\sH,\sH,\sL,\sH)$ gives $p=(2,3,4,5,1,6)$, and interleaving yields
\[
  w \;=\; 2\,3\,4\,5\,2\,3\,1\,4\,5\,6\,1\,6 .
\]
One checks directly from the definitions that $w$ is a nonnesting avoider
of $\{1132,3312\}$ and that $\Phi$ maps it back to $\sC\sB\sC\sC\sA\sC$.
\end{example}

\begin{remark}
The encoding explains the shape of the formula bijectively: each position
of $\sigma$ records a three-way choice --- ``new low''
($\sA$), ``new high'' ($\sB$), or ``sign forced by the geometry'' ($\sC$)
--- for an associated arc, with the position-to-arc correspondence of
Definition~\ref{def:phi} (position $j$ describes arc $j$, except that in
case (I) the first-run letters are shifted one slot left), and with
position $1$ absorbing the case distinction in cases (S) and (II). The two excluded
families are exactly the codes that would describe nothing: a $\sC$-free
code must describe the staircase, whose first letter carries no information
(making the choice $\sB$ redundant: $X_1$), and a leading $\sC$ promises a
gap opener, which forces a nonempty constant block \emph{and} a later free
arc; the codes $\sC v\sC^k$ lack one or the other ($X_2$).
\end{remark}

\appendix

\section{Machine verification}\label{app:verification}

The conjecture was checked by its authors for $n\le8$
\cite[Table~4]{ElizaldeLuo}. Independently of the proofs in this paper,
every lemma above has been verified mechanically, in most cases far beyond
that range, by multiple independently written programs (Python and Rust,
plus clean-room re-implementations written by a different model family
during adversarial review; see Appendix~\ref{app:methods}). Ground truth
for $n\le8$ was established by three independent enumerators working from
the literal containment definition of Section~\ref{sec:conventions} and
cross-validated against each other, giving
\[
 (c_n)_{n=1}^{8} = 1,\,4,\,16,\,58,\,196,\,634,\,1996,\,6178,
\]
in exact agreement with $3^n-3\cdot2^{n-1}+1$, and nonnesting totals
$n!\,\Cat_n=1$, $4$, $30$, $336$, $5040$, $95040$, $2162160$, $57657600$.
Ground-truth avoider
\emph{lists} for $n=9,10$ ($18{,}916$ and $57{,}514$ words --- two levels
beyond the conjecture's original range) were generated by a pruned
depth-first search that is independent of the shape classification of
Section~\ref{sec:shapes}. The table below summarizes the principal checks.

\begin{center}
\small
\begin{tabular}{p{0.42\textwidth}p{0.33\textwidth}p{0.16\textwidth}}
\toprule
\textbf{Claim checked} & \textbf{Scope} & \textbf{Result}\\
\midrule
Fast nonnesting/avoidance tests $\equiv$ literal biconditional containment
 & all $113{,}400$ words of $[5]_2$ and all words for $n\le4$
 & $0$ mismatches\\[2pt]
Theorem~\ref{thm:A}, word level: predicted $\iff$ actual avoider, every
(shape, permutation) pair
 & all $n\le7$; $2{,}162{,}160$ pairs at $n=7$
 & $0$ mismatches\\[2pt]
Theorem~\ref{thm:A} at $n=8$, word level, full exhaustive
 & all $57{,}657{,}600$ pairs (clean-room Rust harness written by the
   adversarial referee)
 & $0$ mismatches; $6{,}178$ avoiders both ways\\[2pt]
Randomized hunts beyond the exhaustive range
 & $n=9,10,11$: $2{,}000{,}000$ uniform (shape, permutation) samples each,
   plus $2{,}000{,}000$ (shape, prefix-interval permutation) samples each
 & $0$ mismatches\\[2pt]
Classification (Lemmas~\ref{lem:obstructions}--\ref{lem:classification}):
$(C)$ satisfiable $\iff$ admissible; per-shape count $=2^{|F(s)|}$
 & brute force over all $2^{n-1}$ sign words per shape, $n\le9$; structural
   checks for every Dyck shape $n\le13$ ($742{,}900$ shapes at $n=13$)
 & exact equality\\[2pt]
Per-shape counts vs.\ independent Rust per-shape ground truth
 & all $1{,}430$ shapes at $n=8$
 & exact match\\[2pt]
$\sum_s N(s)$ over all Dyck shapes, with $N(s)=0$ for inadmissible $s$,
$=3^n-3\cdot2^{n-1}+1$
 & $n\le12$ ($525{,}298$ at $n=12$, matching A168583)
 & pass\\[2pt]
Summation identity of Theorem~\ref{thm:main} as exact integer identity
 & $n\le60$
 & pass\\[2pt]
$|\Wn|$ and Definition~\ref{def:Wn} $\equiv$ complement description
(Lemma~\ref{lem:Wcount})
 & full $3^n$ cube, $n\le10$
 & pass\\[2pt]
Bijection $\Phi$: well-defined, injective, image $=\Wn$,
$\Psi\circ\Phi=\mathrm{id}$, $\Phi\circ\Psi=\mathrm{id}$
 & every valid pair $n\le8$; end-to-end on all $85{,}513$ ground-truth
   avoiders for $n\le10$ (incl.\ the independent DFS lists at $n=9,10$)
 & pass\\
\bottomrule
\end{tabular}
\end{center}

In addition, the proofs went through six independent adversarial audits
(three for the route of Sections~\ref{sec:shapes}--\ref{sec:count}, three
for the bijection of Section~\ref{sec:bijection}), each combining
line-by-line hand verification of every inference with a fresh clean-room
verification suite, and each including a hostile-referee round by a model
from a different family with full code-execution access; all six returned
the verdict \emph{sound}, with no mathematical gap found. The worked
examples (Examples~\ref{ex:n3}, \ref{ex:table} and~\ref{ex:cbccac}) were
re-derived directly from the literal containment definition.

The small-$n$ layer of this table can be reproduced from scratch with the
self-contained script \texttt{\_recompute.py} that accompanies this paper:
it implements only the literal definitions and the statements of the lemmas
(no code is shared with the provers or the earlier verification suites) and
re-checks, among other things: the avoider counts for $n\le5$ over all
words of $\nn$; Theorem~\ref{thm:A} for all pairs with $n\le6$ in the
default run, and with $n\le7$ under the optional flag \texttt{-{}-n7}; the
per-shape counts $N(s)=2^{|F(s)|}$ by brute force over all sign words for
$n\le9$, and the classification totals for $n\le12$; the language $\Wn$ for
$n\le10$; the bijection for $n\le8$; and all three worked examples
(\texttt{ALL CHECKS PASS}; on the preparation machine, under CPython~3.11,
about fifteen seconds for the default suite and a few seconds for the $n=7$
sweep). The heavier artifacts quoted in the table above --- the clean-room
Rust harness used for the exhaustive $n=8$ sweep (recompiled from source
and re-run during the final referee pass: $57{,}657{,}600$ pairs, $0$
mismatches), the shape-level scripts run to $n=13$, the DFS-generated
avoider lists for $n=9,10$, and the audit logs --- are preserved in the
verification bundle that accompanies the source of this paper.

\section{Methods disclosure}\label{app:methods}

This paper, including all mathematical content, was produced by an
artificial-intelligence workflow, and we believe readers and referees
should be able to take this fact into account. The conjecture was selected,
attacked, and solved by language-model agents (Anthropic's Claude family)
running in Demonstrandum, an orchestrated, verification-first multi-agent
pipeline: the conventions of
Section~\ref{sec:conventions} were pinned, as verbatim quotations of the
published LaTeX source of \cite{ElizaldeLuo}, in a frozen definitions file
before any proof attempt; six proof drafts
were generated by parallel agents, of which two survived
triage and were merged into the present exposition (the summation endgame
of Section~\ref{sec:count} and the bijection of
Section~\ref{sec:bijection} were found by different agents and are
logically independent); and the drafts were then subjected to the six
adversarial audits described in Appendix~\ref{app:verification}. A model
from a different family (OpenAI's GPT-5.5, accessed through the Codex CLI
with code-execution access) acted throughout as a hostile referee: it was
prompted to find counterexamples and gaps rather than to confirm, it wrote
its own clean-room verification code (including the exhaustive $n=8$ Rust
check of Appendix~\ref{app:verification}), and its objections --- all
presentational --- were incorporated. No human mathematician contributed
mathematical content; the human role was limited to operating the pipeline.

We emphasize the epistemic status of the result: correctness rests on the
proofs printed in this paper, which are elementary, finite, and intended to
be checkable line by line by a human referee in the ordinary way. The
machine verification of Appendix~\ref{app:verification} is corroboration,
not a substitute for the proofs; conversely, no step of the proofs depends
on unpublished computation. Every number quoted in this paper was either
recomputed for this manuscript by the accompanying script, recomputed
during the final hostile-referee passes on the finished manuscript, or
copied from the verification logs of the audits described above. The proofs were
written with eventual formalization in a proof assistant in mind (all
objects are finite words and tuples; all maps are given explicitly in both
directions), but no formalization is claimed here.

\begin{thebibliography}{9}

\bibitem{Archer2019}
K.~Archer, A.~Gregory, B.~Pennington, and S.~Slayden,
Pattern restricted quasi-Stirling permutations,
\emph{Australas. J. Combin.} \textbf{74} (2019), 389--407.

\bibitem{ArcherLaudone}
K.~Archer and R.~P.~Laudone,
Pattern avoidance in non-crossing and non-nesting permutations,
\emph{Enumer. Comb. Appl.} \textbf{6}:1 (2026), Article \#S2R1,
\texttt{doi:10.54550/ECA2026V6S1R1}; arXiv:2502.13309.

\bibitem{Elizalde2021}
S.~Elizalde,
Descents on quasi-Stirling permutations,
\emph{J. Combin. Theory Ser. A} \textbf{180} (2021), Paper No. 105429.

\bibitem{Elizalde2024}
S.~Elizalde,
Descents on nonnesting multipermutations,
\emph{European J. Combin.} \textbf{121} (2024), Paper No. 103846;
arXiv:2204.00165.

\bibitem{ElizaldeLuo}
S.~Elizalde and A.~Luo,
Pattern avoidance in nonnesting permutations,
\emph{Discrete Math. Theor. Comput. Sci.} \textbf{27}:1, Permutation
Patterns 2024, Paper No. 13 (2025), \texttt{doi:10.46298/dmtcs.14885};
arXiv:2412.00336.

\bibitem{GesselStanley1978}
I.~Gessel and R.~P.~Stanley,
Stirling polynomials,
\emph{J. Combin. Theory Ser. A} \textbf{24} (1978), 24--33.

\bibitem{Luo2024}
A.~Luo,
\emph{Pattern Avoidance in Nonnesting Permutations},
undergraduate thesis, Dartmouth College, May 2024.
\url{https://math.dartmouth.edu/theses/undergrad/2024/Luo-thesis.pdf}

\bibitem{OEIS}
OEIS Foundation Inc.,
\emph{The On-Line Encyclopedia of Integer Sequences},
entry A168583, \url{https://oeis.org/A168583} (accessed June 2026).

\bibitem{SimionSchmidt}
R.~Simion and F.~W.~Schmidt,
Restricted permutations,
\emph{European J. Combin.} \textbf{6} (1985), 383--406.

\end{thebibliography}

\end{document}
